\subsection{Measurement reliability of the fast tier}
\label{sec:reliability}
The fast-tier scores above are noisy, so we measured that noise before using them. We bounded it three ways: repeat evaluation on one device, the same entries on other hosts, and the sampling error implied by the item counts.

We re-evaluated the base model and 10 control entries on one GPU, with the same seed, configuration and batch size of 32 (run \texttt{p2-retest-20261104-202946}). All 11 pairs were bit-identical on every reported metric: GSM8K, MATH-500, IFEval, MMLU-Pro, HellaSwag, refusal, over-refusal, KL, WikiText perplexity and loop rate. The mean and maximum absolute difference were 0.000 on each, so repeat evaluation on one device reproduced every score. Because the seed was held fixed, this test says nothing about sampling noise under a different seed. A repeat panel with a second seed would measure that noise, and we have not run one yet.

We then evaluated the same entries on other hosts, under the same seed, configuration and batch size. The local RTX PRO 4500 was compared with a Lambda H100 PCIe on 11 entries and with a Lambda A10 on four of them. Table~\ref{tab:devices} gives the local-versus-H100 differences.
\begin{table}[t]
\centering
\begin{tabular}{lrrrr}\toprule
Benchmark & Items per entry & Mean abs. diff. & Max abs. diff. & Spearman $\rho$ \\\midrule
GSM8K & 50 & 0.053 & 0.12 & 0.30 \\
MATH-500 & 25 & 0.040 & 0.12 & 0.09 \\
IFEval & 25 & 0.022 & 0.08 & 0.29 \\
Over-refusal & 50 & 0.049 & 0.14 & 0.39 \\
MMLU-Pro & 100 & 0.008 & 0.02 & 0.75 \\
HellaSwag & 100 & 0.003 & 0.01 & 0.76 \\
KL to base & 50 & 0.000 & 0.001 & 1.00 \\
Loop rate & 100 responses & 0.006 & 0.04 & 0.72 \\\bottomrule
\end{tabular}
\caption{Fast tier, local RTX PRO 4500 against Lambda H100 PCIe. Same seed, configuration and batch size 32. $n = 11$ entries were evaluated on both hosts. Mean and maximum absolute difference between hosts are taken over entries. Spearman correlation is taken across the 11 entries. The loop rate pools the GSM8K, MATH-500 and IFEval responses of each entry at the 1,024-token fast-tier cap. The entries vary little, so these correlations do not qualify any target. No single run id covers this panel, which was assembled from per-host fast-tier records.}
\label{tab:devices}
\end{table}

In the paragraph below, each pair of values is the local score followed by the other host's. On the base model, the H100 differed from the local host by 0.10 on GSM8K (0.16 and 0.06; 50 items), by 0.08 on MATH-500 (0.16 and 0.08; 25 items) and by 0.10 on over-refusal (0.10 and 0.20; 50 prompts). The other benchmarks agreed within 0.01. Against the A10, the gaps were 0.10 on GSM8K (0.16 and 0.06; 50 items), 0.12 on MATH-500 (0.16 and 0.28; 25 items) and 0.04 on IFEval (0.40 and 0.44; 25 items). MMLU-Pro (0.28) and HellaSwag (0.39) matched on the base model, with 100 items each. On the four edited entries shared with the A10, the largest gaps were 0.16 on MATH-500 and 0.12 on IFEval. MMLU-Pro, HellaSwag, KL and keyword refusal agreed within 0.02 on those entries.

Seed matching does not make sampled generation reproducible across hosts at temperature 0.6. We expect small numerical differences to change a sampled trajectory after its first divergence. Atil et al. report a similar pattern under settings configured to be deterministic \citep{atil2024nondeterminism}.

Item counts set a second floor. For a proportion $p$ on $n$ independent items, the binomial standard error is $\sqrt{p(1-p)/n}$. At the base proportions above, that gives 0.03 to 0.05 for GSM8K ($n = 50$), 0.07 to 0.09 for MATH-500 ($n = 25$), about 0.10 for IFEval ($n = 25$), 0.045 for MMLU-Pro ($n = 100$) and 0.049 for HellaSwag ($n = 100$), or 0.03 to 0.10 across 25 to 100 items. These are planning values, not the final statistic. They assume independent items and ignore host effects. Miller treats evaluation questions as draws from a super-population and gives formulas for evaluation data \citep{miller2024error}. We use the binomial form for single scores.

Those floors set what a claim may say. We therefore report GSM8K, MATH-500 and IFEval as descriptive at the fast tier and make no claim about single-edit effects on them. A correlation against an observed change cannot exceed about the square root of that change's reliability, so each correlation must be reported with the reliability of its target. The H1 threshold of $\rho \ge 0.6$ would need a target reliability of at least 0.36. We have not estimated a reliability coefficient for any fast-tier target, so no correlation in this paper is qualified. Table~\ref{tab:admissible} lists the claims the fast tier supports at present.
\begin{table}[t]
\centering
\begin{tabular}{p{0.36\linewidth}p{0.24\linewidth}p{0.32\linewidth}}\toprule
Claim & Status & Basis \\\midrule
Single-edit change on GSM8K, MATH-500 or IFEval & Not admissible & Binomial SE 0.03 to 0.10; base host gaps up to 0.12 \\
Ranking edits on GSM8K, MATH-500 or IFEval & Not admissible & Inherits the same noise \\
Correlation with an observed change on these three & Only with the target reliability reported & Ceiling of about the square root of reliability; reliability not yet estimated \\
Family-mean difference on these three & Not admissible without host balance or a same-device base & Host bias can shift whole families \\
Keyword refusal change when the base is at 0.99 to 1.00 & Not admissible & At 50 prompts one item is 0.02; the metric sits at its ceiling \\
Loop rate against the blog's MATH-500 share & Not admissible & Token caps (1,024 and 32,768), pooling and sample counts differ \\
MMLU-Pro and HellaSwag family means & Conditional on host-balanced means, a same-device base and intervals & Single-entry SE 0.045 to 0.049 at 100 items \\
WikiText perplexity and KL & Conditional on the numerical floor being stated & Base-model perplexity agreed across hosts to 0.01 \\\bottomrule
\end{tabular}
\caption{Claims supported by the fast tier at present. ``Not admissible'' means no claim of that kind is made. ``Conditional'' means the claim may be made only under the stated conditions. Neither status qualifies a predictor target. Binomial values are planning approximations.}
\label{tab:admissible}
\end{table}

Delta targets for sampled-generation benchmarks must be taken against a base evaluated on the same device. In this panel, MMLU-Pro, HellaSwag and KL agreed closely across hosts, and keyword refusal agreed within 0.02 on the A10 entries. GSM8K, MATH-500, IFEval and over-refusal did not. We have not applied the qualification rule. It requires a Spearman $\rho$ of at least 0.8 with a lower confidence bound of at least 0.6, and this panel has too little variation between entries to test it. MMLU-Pro, HellaSwag, perplexity and KL become predictor targets only after their measurement floors are fixed, and that step needs a second-seed repeat panel.

The D3 labels carry a provenance limit. Runtime metadata for earlier records was backfilled from job logs after the runs. We treat it as inferred, not measured, and the backfill does not borrow runtime fields from another device. Of 300 evaluated D3 edits, 260 had matched-reference delta targets, and all 260 retained labels use inferred runtime metadata (run \texttt{p4/board-a-d3-audit-20261104}). The other 40 had no predecessor target and are excluded. The host of each D3 entry is not yet tabulated, so the D3 results in this paper are not host-balanced.

The fast tier remains useful for descriptive reporting and for effects several times larger than these floors.
